\documentclass[UTF8,11pt]{ctexart}

\usepackage[a4paper,margin=1in]{geometry}
\usepackage{amsmath,amssymb,amsthm,mathtools}
\usepackage{bm}
\usepackage{bbm}
\usepackage{enumitem}
\usepackage{hyperref}
\usepackage{tikz-cd}

\title{Coefficient-Split 视角下 CKKS Bootstrapping Core 的子环分解}
\author{}
\date{}

% theorem environments
\newtheorem{theorem}{定理}
\newtheorem{lemma}{引理}
\newtheorem{proposition}{命题}
\newtheorem{corollary}{推论}
\theoremstyle{definition}
\newtheorem{definition}{定义}
\newtheorem{remark}{注记}
\newtheorem{example}{例子}

% commands
\newcommand{\C}{\mathbb{C}}
\newcommand{\Z}{\mathbb{Z}}
\newcommand{\RN}{R_N}
\newcommand{\Rn}{R_n}
\newcommand{\Eval}{\operatorname{Eval}}
\newcommand{\CtwoS}{\operatorname{C2S}}
\newcommand{\StwoC}{\operatorname{S2C}}
\newcommand{\BTS}{\operatorname{BTS}}
\newcommand{\Split}{\operatorname{Split}}
\newcommand{\Merge}{\operatorname{Merge}}
\newcommand{\EvalMod}{\operatorname{EvalMod}}
\newcommand{\Diag}{\operatorname{diag}}
\newcommand{\Id}{\operatorname{Id}}
\newcommand{\BlockSplit}{\operatorname{BlockSplit}}
\newcommand{\Red}{\operatorname{Red}}
\newcommand{\op}{\operatorname}

\begin{document}

\maketitle

\begin{abstract}
本文研究 CKKS bootstrapping core
\[
\CtwoS_N \longrightarrow \EvalMod_N \longrightarrow \StwoC_N
\]
是否可以沿 coefficient-domain subring split 分解为多个小环上的 leaf bootstrapping core。
我们关注的分解形式是
\[
\Split
\longrightarrow
\bigoplus_{r=0}^{k-1}\CtwoS_n
\longrightarrow
\bigoplus_{r=0}^{k-1}\EvalMod_n
\longrightarrow
\bigoplus_{r=0}^{k-1}\StwoC_n
\longrightarrow
\Merge,
\]
其中 \(N=kn\)。核心观点是：当前路线使用的是 coefficient split，而不是 slot-preserving split。
虽然 coefficient split 在 slot 视角下会把大环 slot 做成局部线性组合，但它与 leaf C2S 组合后，恰好等价于先做大环 C2S 再按 coefficient residue class 分组。
因此，如果 EvalMod 是同一个逐坐标标量函数，则 leaf EvalMod 可以独立作用于每个 coefficient 子序列，从而得到 no-\(B_k\) 的 subring bootstrapping decomposition。
\end{abstract}

\tableofcontents

\section{研究问题}

标准 CKKS bootstrapping 通常可以分成如下几步：
\[
\operatorname{ScaleDown}
\longrightarrow
\operatorname{ModUp}
\longrightarrow
\CtwoS_N
\longrightarrow
\EvalMod_N
\longrightarrow
\StwoC_N.
\]
本文暂时聚焦其中的 core：
\[
\BTS_N
:=
\StwoC_N\circ \EvalMod_N\circ \CtwoS_N.
\]

我们希望理解如下分解是否在理论上成立：
\[
\BTS_N
\stackrel{?}{\approx}
\gamma_k^{-1}
\Merge
\circ
\left(
\bigoplus_{r=0}^{k-1}\BTS_n
\right)
\circ
\Split,
\]
其中 \(N=kn\)，\(\gamma_k\) 是可能由实现 convention 引入的全局 normalization factor。

等价地，我们想研究大环 bootstrapping core
\[
\CtwoS_N
\longrightarrow
\EvalMod_N
\longrightarrow
\StwoC_N
\]
是否可以替换为
\[
\Split
\longrightarrow
\bigoplus_{r=0}^{k-1}\CtwoS_n
\longrightarrow
\bigoplus_{r=0}^{k-1}\EvalMod_n
\longrightarrow
\bigoplus_{r=0}^{k-1}\StwoC_n
\longrightarrow
\Merge.
\]

这个问题并不显然，因为：
\begin{enumerate}[label=(\arabic*)]
    \item \(\Split\) 是 coefficient-domain split；
    \item CKKS bootstrapping core 中的 \(\CtwoS\)、\(\EvalMod\)、\(\StwoC\) 是通过 slot/canonical embedding 实现；
    \item coefficient split 后的小环 slots 并不是大环 slots 的简单子集；
    \item \(\EvalMod\) 是非线性操作，一般不能任意穿过线性变换。
\end{enumerate}

因此本文的理论路线是：
\begin{enumerate}[label=(\roman*)]
    \item 先澄清 coefficient split 在 slot 视角下的语义；
    \item 证明 \(\Split\) 与 \(\CtwoS\) 的交换图；
    \item 分析 \(\EvalMod\) 作为逐坐标标量函数时与 coefficient 分组的交换性；
    \item 合成 no-\(B_k\) 的 bootstrapping decomposition 主定理；
    \item 最后解释 \(B_k\) 的含义，以及为什么当前路线不需要显式 \(B_k\)。
\end{enumerate}

\section{大环、小环与 canonical embedding}

设
\[
R_N=\C[X]/(X^N+1),
\]
其中 \(N=kn\)。定义小环
\[
R_n=\C[Y]/(Y^n+1),
\]
并令
\[
Y=X^k.
\]

大环的 roots 满足
\[
\alpha^N=-1.
\]
小环的 roots 满足
\[
\theta^n=-1.
\]
如果 \(\alpha^N=-1\)，则
\[
(\alpha^k)^n=\alpha^{kn}=\alpha^N=-1,
\]
所以
\[
\theta=\alpha^k
\]
是小环的 root。

设
\[
R_N=\mathbb{C}[X]/(X^N+1).
\]
令 \(\alpha_0,\ldots,\alpha_{N-1}\) 为 \(X^N+1\) 的全部复根，即
\[
\alpha_j^N+1=0,
\]
等价地，
\[
\alpha_j^N=-1.
\]

由于如果 \(P(X)\equiv Q(X)\pmod{X^N+1}\)，则对任意根 \(\alpha_j\) 都有
\[
P(\alpha_j)=Q(\alpha_j),
\]
所以在这些根上求值给出了商环上的良定义映射。该映射称为 canonical embedding：
\[
\sigma_N:R_N\to \mathbb{C}^N,
\qquad
\sigma_N(P)
=
\bigl(P(\alpha_0),P(\alpha_1),\ldots,P(\alpha_{N-1})\bigr).
\]

也可以简写为
\[
\sigma_N(P)=\left(P(\alpha)\right)_{\alpha^N=-1}.
\]

在 CKKS 中，由于共轭对坐标冗余，实际可独立编码的 complex slots 通常为 \(N/2\) 个；但为了推导线性关系，本文先使用 full canonical embedding。

类似地，对于小环多项式
\[
Q(Y)=\sum_{j=0}^{n-1}q_jY^j,
\]
其小环 slot vector 为
\[
\sigma_n(Q)
=
\left(Q(\theta)\right)_{\theta^n=-1}.
\]

\begin{remark}
真实 CKKS 中通常只使用 conjugate-invariant 的一半 complex slots，并且实现中还会出现 bit-reversal、slot permutation、real/imag split、scale convention 等工程细节。
本文先在 full complex embedding 的理想模型下推导核心结构。真实实现中的 permutation 和 scale 可以视为在该理想图上额外复合的可逆线性 convention。
\end{remark}

\section{Coefficient split 与 slot-preserving split}

本节说明两种不同的“拆分”：

\begin{enumerate}[label=(\arabic*)]
    \item coefficient split：按多项式系数的 residue class 拆分；
    \item slot-preserving split：希望小环 slots 原样保存某些大环 slots。
\end{enumerate}

这两者不是同一个操作。当前 no-\(B_k\) 路线使用的是 coefficient split。

\subsection{\(k=2\) 的 coefficient split}

先考虑 \(k=2\)，即 \(N=2n\)。任意
\[
P(X)=a_0+a_1X+a_2X^2+\cdots+a_{N-1}X^{N-1}
\]
都可以唯一写为
\[
P(X)=A(X^2)+XB(X^2),
\]
其中
\[
A(Y)=a_0+a_2Y+a_4Y^2+\cdots,
\]
\[
B(Y)=a_1+a_3Y+a_5Y^2+\cdots.
\]

用系数向量表示，令
\[
p=(a_0,a_1,a_2,a_3,\ldots,a_{N-1})^T.
\]
定义偶数系数选择矩阵 \(E\) 与奇数系数选择矩阵 \(O\)：
\[
Ep=(a_0,a_2,a_4,\ldots)^T,
\]
\[
Op=(a_1,a_3,a_5,\ldots)^T.
\]
则 coefficient split 是
\[
p\longmapsto (Ep,Op).
\]

\subsection{Coefficient split 在 slot 视角下的作用}

令 \(\alpha\) 是大环 root，\(\alpha^N=-1\)。当 \(N=2n\) 时，令
\[
\theta=\alpha^2.
\]
则
\[
\theta^n=(\alpha^2)^n=\alpha^N=-1,
\]
所以 \(\theta\) 是小环 root。

大环的一对 roots 可以取为
\[
\alpha,\quad -\alpha,
\]
它们都映射到同一个小环 root：
\[
\alpha^2=(-\alpha)^2=\theta.
\]

由
\[
P(X)=A(X^2)+XB(X^2)
\]
可得
\[
P(\alpha)=A(\theta)+\alpha B(\theta),
\]
以及
\[
P(-\alpha)=A(\theta)-\alpha B(\theta).
\]

记
\[
z^+=P(\alpha),\qquad z^-=P(-\alpha),
\]
\[
u=A(\theta),\qquad v=B(\theta).
\]
则
\[
z^+=u+\alpha v,
\]
\[
z^-=u-\alpha v.
\]
反过来，
\[
u=\frac{z^++z^-}{2},
\]
\[
v=\frac{z^+-z^-}{2\alpha}.
\]

因此 coefficient split 在 slot 视角下不是简单保留大环 slots，而是
\[
(P(\alpha),P(-\alpha))
\longmapsto
\left(
\frac{P(\alpha)+P(-\alpha)}{2},
\frac{P(\alpha)-P(-\alpha)}{2\alpha}
\right).
\]

\begin{remark}
这说明：coefficient split 后的小环 slots
\[
A(\theta),\quad B(\theta)
\]
不是原始大环 slots
\[
P(\alpha),\quad P(-\alpha)
\]
的子集，而是它们的局部线性组合。
\end{remark}

\subsection{Slot-preserving split}

如果目标不是 coefficient split，而是希望构造两个小环多项式
\[
Q^+(Y),\quad Q^-(Y)
\]
使得
\[
Q^+(\theta)=P(\alpha),
\]
\[
Q^-(\theta)=P(-\alpha),
\]
那么仅仅输出 \(A(Y),B(Y)\) 不够。

由前面的关系，
\[
P(\alpha)=A(\theta)+\alpha B(\theta).
\]
如果存在小环多项式 \(h(Y)\)，满足
\[
h(\theta)=\alpha,\qquad \alpha^2=\theta,
\]
那么可定义
\[
Q^+(Y)=A(Y)+h(Y)B(Y),
\]
\[
Q^-(Y)=A(Y)-h(Y)B(Y).
\]
于是
\[
Q^+(\theta)
=
A(\theta)+h(\theta)B(\theta)
=
A(\theta)+\alpha B(\theta)
=
P(\alpha),
\]
同理
\[
Q^-(\theta)=P(-\alpha).
\]

由于
\[
h(\theta)^2=\alpha^2=\theta,
\]
所以在商环意义下
\[
h(Y)^2\equiv Y\pmod{Y^n+1}.
\]
因此 \(h(Y)\) 可以理解为小环上 \(Y\) 的某个平方根选择多项式。
\subsection{Slot-preserving split 的矩阵形式}

上一节中我们已经看到，普通的 coefficient split 给出
\[
P(X)=A(X^2)+XB(X^2),
\]
其中
\[
A(Y)=a_0+a_2Y+a_4Y^2+\cdots,
\]
\[
B(Y)=a_1+a_3Y+a_5Y^2+\cdots.
\]

也就是说，如果
\[
p=(a_0,a_1,a_2,a_3,\ldots,a_{N-1})^T
\]
是 \(P(X)\) 的系数向量，那么 \(A(Y)\) 和 \(B(Y)\) 的系数向量分别是
\[
Ep=(a_0,a_2,a_4,\ldots)^T,
\]
\[
Op=(a_1,a_3,a_5,\ldots)^T.
\]

这里 \(E\) 表示取偶数下标系数的选择矩阵，\(O\) 表示取奇数下标系数的选择矩阵。

因此，coefficient split 在系数向量上的形式就是
\[
p\longmapsto (Ep,Op).
\]

但是 coefficient split 并不是 slot-preserving split。它输出的小环多项式是
\[
A(Y),\quad B(Y),
\]
它们的小环 slots 是
\[
A(\theta_j),\quad B(\theta_j).
\]

而如果我们希望两个小环多项式直接保存原来的大环 slot pair：
\[
P(\alpha_j),\quad P(-\alpha_j),
\]
那么需要构造两个新的小环多项式
\[
Q^+(Y),\quad Q^-(Y),
\]
满足
\[
Q^+(\theta_j)=P(\alpha_j),
\]
\[
Q^-(\theta_j)=P(-\alpha_j).
\]

由前面的关系
\[
P(\alpha_j)=A(\theta_j)+\alpha_jB(\theta_j),
\]
\[
P(-\alpha_j)=A(\theta_j)-\alpha_jB(\theta_j),
\]
我们知道，只要存在一个小环多项式 \(h(Y)\)，使得
\[
h(\theta_j)=\alpha_j,
\]
就可以定义
\[
Q^+(Y)=A(Y)+h(Y)B(Y),
\]
\[
Q^-(Y)=A(Y)-h(Y)B(Y).
\]

这样就有
\[
Q^+(\theta_j)
=
A(\theta_j)+h(\theta_j)B(\theta_j)
=
A(\theta_j)+\alpha_jB(\theta_j)
=
P(\alpha_j),
\]
以及
\[
Q^-(\theta_j)
=
A(\theta_j)-h(\theta_j)B(\theta_j)
=
A(\theta_j)-\alpha_jB(\theta_j)
=
P(-\alpha_j).
\]

下面我们把这个构造写成矩阵形式。

设小环 \(R_n=\mathbb{C}[Y]/(Y^n+1)\) 的 roots 为
\[
\theta_0,\theta_1,\ldots,\theta_{n-1}.
\]

定义小环 evaluation matrix \(F_n\)。如果
\[
q=(q_0,q_1,\ldots,q_{n-1})^T
\]
是某个小环多项式
\[
Q(Y)=q_0+q_1Y+\cdots+q_{n-1}Y^{n-1}
\]
的系数向量，那么
\[
F_n q
=
\begin{pmatrix}
Q(\theta_0)\\
Q(\theta_1)\\
\vdots\\
Q(\theta_{n-1})
\end{pmatrix}.
\]

换句话说，\(F_n\) 把小环多项式的系数向量变成它的小环 slot 向量。

现在令
\[
D=\operatorname{Diag}(\alpha_0,\alpha_1,\ldots,\alpha_{n-1}).
\]

这个矩阵 \(D\) 是 slot 坐标中的逐坐标乘法矩阵。也就是说，如果
\[
v=
\begin{pmatrix}
B(\theta_0)\\
B(\theta_1)\\
\vdots\\
B(\theta_{n-1})
\end{pmatrix}
\]
是 \(B(Y)\) 的 slot 向量，那么
\[
Dv
=
\begin{pmatrix}
\alpha_0B(\theta_0)\\
\alpha_1B(\theta_1)\\
\vdots\\
\alpha_{n-1}B(\theta_{n-1})
\end{pmatrix}.
\]

由于 \(h(\theta_j)=\alpha_j\)，所以在 slot 坐标中，乘以多项式 \(h(Y)\) 的效果正是
\[
B(\theta_j)\longmapsto h(\theta_j)B(\theta_j)=\alpha_jB(\theta_j).
\]

因此，乘以 \(h(Y)\) 在 slot 坐标中对应矩阵 \(D\)。

但是我们现在想要的是它在系数坐标中的矩阵形式。设 \(b\) 是 \(B(Y)\) 的系数向量：
\[
b=Op.
\]
那么 \(B(Y)\) 的 slot 向量是
\[
F_n b.
\]

在 slot 坐标中乘以 \(h(Y)\)，得到
\[
D F_n b.
\]

再把结果从 slot 坐标变回系数坐标，需要乘以 \(F_n^{-1}\)。因此，系数坐标中的操作是
\[
b\longmapsto F_n^{-1}D F_n b.
\]

我们定义
\[
H:=F_n^{-1}DF_n.
\]

于是 \(H\) 就是“乘以 \(h(Y)\)”这个操作在小环系数坐标中的矩阵。

也就是说，如果 \(b=Op\) 是 \(B(Y)\) 的系数向量，那么
\[
HOp
\]
就是 \(h(Y)B(Y)\) 的系数向量。

另一方面，\(A(Y)\) 的系数向量是
\[
Ep.
\]

因此
\[
Q^+(Y)=A(Y)+h(Y)B(Y)
\]
的系数向量是
\[
q^+=Ep+HOp.
\]

同理，
\[
Q^-(Y)=A(Y)-h(Y)B(Y)
\]
的系数向量是
\[
q^-=Ep-HOp.
\]

所以 slot-preserving split 在系数向量上的形式为
\[
p\longmapsto(q^+,q^-)
=
(Ep+HOp,\;Ep-HOp).
\]

这和 coefficient split
\[
p\longmapsto(Ep,Op)
\]
完全不同。

\begin{remark}
Coefficient split 只是把 \(P(X)\) 的偶数系数和奇数系数分开：
\[
p\mapsto(Ep,Op).
\]

而 slot-preserving split 的目标是让两个小环多项式的 slots 分别等于原来的大环 slot pair：
\[
Q^+(\theta_j)=P(\alpha_j),
\qquad
Q^-(\theta_j)=P(-\alpha_j).
\]

为了做到这一点，需要额外把 \(B(Y)\) 乘以平方根选择多项式 \(h(Y)\)。这个操作在系数坐标中不是简单的取系数，而是一个线性变换
\[
H=F_n^{-1}DF_n.
\]

因此 slot-preserving split 是
\[
p\mapsto(Ep+HOp,\;Ep-HOp),
\]
而不是
\[
p\mapsto(Ep,Op).
\]

当前 no-\(B_k\) bootstrapping decomposition 使用的是 coefficient split，也就是
\[
P(X)\mapsto(A(Y),B(Y)),
\]
并不试图让 leaf slots 原样保存
\[
P(\alpha_j),P(-\alpha_j).
\]
这也是为什么当前路线不需要先构造 \(Q^+,Q^-\)，也不需要显式乘以 \(h(Y)\)。
\end{remark}

\section{Split 与 C2S 的交换图}

本节证明 no-\(B_k\) 路线最核心的线性恒等式：
\[
(\CtwoS_n\oplus \CtwoS_n)\circ \Split
=
(E/O)\circ \CtwoS_N.
\]

直观地说，这个等式表示：

\begin{quote}
先做 coefficient split，再分别对两个 leaf 做 C2S，等价于先对大环做 C2S，再把得到的 coefficient slots 按偶数/奇数分组。
\end{quote}

\subsection{Evaluation matrix 表示}

令 \(F_N\) 是大环 evaluation matrix：
\[
F_Np=\sigma_N(P).
\]
理想情况下，
\[
\CtwoS_N=F_N^{-1},
\qquad
\StwoC_N=F_N.
\]

类似地，小环有
\[
\CtwoS_n=F_n^{-1},
\qquad
\StwoC_n=F_n.
\]

\begin{remark}
    C2S 叫 CoeffsToSlots。 这个名字的意思不是： $p \mapsto F_N p$。
    
    而是：
    把明文多项式的 coefficients p 放到 CKKS slots 里面，让后面的 EvalMod 能逐 slot 作用在这些 coefficients 上。
\end{remark}

\subsection{\(k=2\) 情况的证明}

输入是大环 slot vector
\[
z=F_Np.
\]
Coefficient split 后得到两个小环多项式 \(A,B\)，其系数分别是
\[
Ep,\qquad Op.
\]
它们的小环 slot vector 是
\[
F_nEp,\qquad F_nOp.
\]

因此，\(\Split\) 在 slot 视角下的矩阵为
\[
\Split^{slot}
=
\begin{pmatrix}
F_nE\\
F_nO
\end{pmatrix}
F_N^{-1}.
\]

现在对两个 leaf 分别做 \(\CtwoS_n\)：
\[
\CtwoS_n\oplus\CtwoS_n
=
\begin{pmatrix}
F_n^{-1} & 0\\
0 & F_n^{-1}
\end{pmatrix}.
\]
于是
\[
(\CtwoS_n\oplus\CtwoS_n)\circ\Split^{slot}
=
\begin{pmatrix}
F_n^{-1} & 0\\
0 & F_n^{-1}
\end{pmatrix}
\begin{pmatrix}
F_nE\\
F_nO
\end{pmatrix}
F_N^{-1}.
\]
化简得
\[
=
\begin{pmatrix}
E\\
O
\end{pmatrix}
F_N^{-1}.
\]
而
\[
(E/O)\circ\CtwoS_N
=
\begin{pmatrix}
E\\
O
\end{pmatrix}
F_N^{-1}.
\]
因此
\[
\boxed{
(\CtwoS_n\oplus\CtwoS_n)\circ\Split
=
(E/O)\circ\CtwoS_N.
}
\]

\begin{remark}
这个等式解释了为什么 no-\(B_k\) 路线在坐标上是自洽的。
虽然 \(\Split\) 后的小环 slots 不是原始大环 slots 的子集，但再接 leaf C2S 后，得到的正是大环 C2S 输出的 coefficient 子序列。
\end{remark}

\subsection{一般 \(k\) 的推广}

令 \(N=kn\)。任意大环多项式可以写为
\[
P(X)=\sum_{r=0}^{k-1}X^rP_r(X^k),
\]
其中
\[
P_r(Y)=\sum_{j}a_{r+jk}Y^j.
\]

定义 residue class coefficient selection：
\[
P_r^{\mathrm{sel}}p
=
(a_r,a_{r+k},a_{r+2k},\ldots)^T.
\]
为避免和多项式 \(P_r\) 混淆，下文把选择矩阵记为 \(S_r\)：
\[
S_rp=(a_r,a_{r+k},a_{r+2k},\ldots)^T.
\]

Coefficient split 为
\[
p\longmapsto (S_0p,S_1p,\ldots,S_{k-1}p).
\]
在 slot 视角下：
\[
\Split_k^{slot}
=
\begin{pmatrix}
F_nS_0\\
F_nS_1\\
\vdots\\
F_nS_{k-1}
\end{pmatrix}
F_N^{-1}.
\]

对每个 leaf 做 \(\CtwoS_n=F_n^{-1}\)，得到：
\[
\left(\bigoplus_{r=0}^{k-1}\CtwoS_n\right)\circ\Split_k
=
\begin{pmatrix}
S_0\\
S_1\\
\vdots\\
S_{k-1}
\end{pmatrix}
F_N^{-1}.
\]
而右边就是先做大环 C2S，再按 residue class 分组：
\[
\left(\bigoplus_{r=0}^{k-1}S_r\right)\circ\CtwoS_N
=
\begin{pmatrix}
S_0\\
S_1\\
\vdots\\
S_{k-1}
\end{pmatrix}
F_N^{-1}.
\]

因此有一般形式：

\[
\boxed{
\left(\bigoplus_{r=0}^{k-1}\CtwoS_n\right)\circ\Split_k
=
\left(\bigoplus_{r=0}^{k-1}S_r\right)\circ\CtwoS_N.
}
\]

这就是 no-\(B_k\) 路线的第一个核心交换图。

\section{EvalMod 与 coefficient 分组的交换性}

上一节处理了线性层：
\[
\CtwoS.
\]
本节处理唯一的非线性层：
\[
\EvalMod.
\]

\subsection{EvalMod 的抽象模型}

在 CKKS bootstrapping 中，经过 \(\CtwoS\) 之后，slots 中装入的是 coefficient-like values。
更具体地说，\(\CtwoS_N\) 的作用可以抽象理解为：把大环明文多项式的系数（注意，下面是$a$不是$\alpha$）
\[
(a_0,a_1,\ldots,a_{N-1})
\]
搬到 slots 中，使得后续的 \(\EvalMod\) 可以逐坐标地作用在这些系数型数据上。

在 ModUp 之后，这些 coefficient-like values 可以抽象为
\[
x_i=m_i+q_0t_i+\varepsilon_i,
\]
其中：
\begin{itemize}
    \item \(m_i\) 是我们希望恢复的明文系数；
    \item \(q_0\) 是底层 modulus；
    \item \(t_i\in\mathbb Z\) 是 ModUp 后产生的整数 wrap 项；
    \item \(\varepsilon_i\) 表示噪声或近似误差。
\end{itemize}

忽略噪声时，可以写成
\[
x_i=m_i+q_0t_i.
\]

$\EvalMod$ 的目标就是逐坐标地去掉 \(q_0t_i\)，恢复 \(m_i\)：
\[
x_i=m_i+q_0t_i
\longmapsto
m_i.
\]

等价地，如果进行 mod1 归一化，令
\[
y_i=\frac{x_i}{q_0}
=
t_i+\frac{m_i}{q_0}.
\]
记
\[
\epsilon_i=\frac{m_i}{q_0}.
\]
那么
\[
y_i=t_i+\epsilon_i,
\qquad t_i\in\mathbb Z.
\]
此时 \(\EvalMod\) 的任务可以理解为从
\[
t_i+\epsilon_i
\]
中恢复小量
\[
\epsilon_i.
\]
也就是说，它近似实现
\[
y_i \bmod 1.
\]

这也是 CKKS bootstrapping 中 EvalMod 常常与 scaled sine 或 cosine 近似有关的原因。
因为
\[
\sin(2\pi y_i)
=
\sin(2\pi(t_i+\epsilon_i))
=
\sin(2\pi\epsilon_i)
\approx
2\pi\epsilon_i
\]
当 \(\epsilon_i\) 足够小时成立。
因此，适当缩放的 sine 函数可以近似恢复 \(\epsilon_i\)。
实际实现中通常不会直接计算 sine，而是使用某种多项式近似，例如 Chebyshev 近似、minimax 近似或结合 double-angle 技巧的近似电路。

在本文的抽象模型中，我们用一个标量函数
\[
f
\]
表示 EvalMod 在单个 coordinate 上实现的近似模约减。它满足
\[
f(m_i+q_0t_i)\approx m_i,
\]
或者在归一化形式下满足
\[
f(t_i+\epsilon_i)\approx \epsilon_i.
\]

由于 EvalMod 是逐坐标作用的，我们将它在 \(N\) 个坐标上的作用记为
\[
f^{\oplus N}.
\]

也就是说，对向量
\[
x=(x_0,x_1,\ldots,x_{N-1})
\]
定义
\[
f^{\oplus N}(x)
=
(f(x_0),f(x_1),\ldots,f(x_{N-1})).
\]

因此，在理想抽象模型中，大环 EvalMod 可以写成
\[
\EvalMod_N=f^{\oplus N}.
\]

类似地，小环上有 \(n\) 个坐标，leaf EvalMod 可以写成
\[
\EvalMod_n=f^{\oplus n}.
\]

这里需要注意，\(f^{\oplus N}\) 一般不是线性映射，也不是矩阵；它只是同一个标量函数 \(f\) 在每个坐标上独立作用得到的向量函数。
这个记号也避免了与前文中的 evaluation matrix \(F_N\) 混淆。

\subsection{Residue-class 分组算子}

为了描述 coefficient split 后的各个 leaf，我们需要定义一个按 residue class 选择系数的分组算子。

设
\[
N=kn.
\]
令
\[
x=(x_0,x_1,\ldots,x_{N-1})^T
\]
是一个长度为 \(N\) 的 coefficient-packed vector。

\begin{remark}
    这里需要区分两个向量。设 ModUp 后的明文多项式为
\[
P(X)=p_0+p_1X+\cdots+p_{N-1}X^{N-1}.
\]
记它的系数向量为
\[
p=(p_0,p_1,\ldots,p_{N-1})^T.
\]
它的 canonical slot vector 是
\[
z=\sigma_N(P)=F_Np.
\]

在 bootstrapping core 中，\(\CtwoS_N\) 的作用是把当前 canonical slots
\[
z=F_Np
\]
变成装有 coefficients 的 slots：
\[
\CtwoS_N(z)=p.
\]

因此，下面我们把
\[
x=(x_0,x_1,\ldots,x_{N-1})^T
\]
称为 C2S 后的 coefficient-packed slot vector，意思是：
\[
x=p.
\]
也就是说，\(x_i\) 是 C2S 前多项式 \(P(X)\) 的第 \(i\) 个系数，但它现在作为 C2S 后密文的 slot 内容出现。

注意，\(x\) 不是 \(P(X)\) 的 canonical embedding；\(P(X)\) 的 canonical embedding 是
\[
\sigma_N(P)=F_Np.
\]

在 ModUp 后，这些系数型 slot 值可以抽象为
\[
x_i=m_i+q_0t_i+\varepsilon_i.
\]
\end{remark}

对于每个
\[
r=0,1,\ldots,k-1,
\]
定义选择算子
\[
S_r:\mathbb{C}^N\to \mathbb{C}^n
\]
为
\[
S_rx
=
(x_r,x_{r+k},x_{r+2k},\ldots,x_{r+(n-1)k})^T.
\]

也就是说，\(S_r\) 从大环 coefficient-packed vector 中取出所有下标与 \(r\) 同余的坐标。

例如，当 \(k=2\) 时，
\[
S_0x=(x_0,x_2,x_4,\ldots)^T,
\]
\[
S_1x=(x_1,x_3,x_5,\ldots)^T.
\]

这正对应于 even/odd coefficient split。

更一般地，我们把所有 residue class 的选择算子合在一起，定义整体分组算子
\[
S:\mathbb{C}^N\to (\mathbb{C}^n)^k
\]
为
\[
Sx
=
(S_0x,S_1x,\ldots,S_{k-1}x).
\]

矩阵形式下，可以写为
\[
S=
\begin{pmatrix}
S_0\\
S_1\\
\vdots\\
S_{k-1}
\end{pmatrix}.
\]

因此
\[
Sx=
\begin{pmatrix}
S_0x\\
S_1x\\
\vdots\\
S_{k-1}x
\end{pmatrix}.
\]

如果
\[
x=(x_0,x_1,\ldots,x_{N-1})^T
\]
是大环 C2S 后得到的 coefficient-packed vector，那么
\[
S_rx=(x_r,x_{r+k},x_{r+2k},\ldots)^T
\]
就是第 \(r\) 个 leaf 在 leaf C2S 后应该得到的 coefficient-packed vector。

换句话说，\(S\) 描述的是：
\[
\text{大环 coefficient-packed vector}
\longmapsto
\text{按 residue class 分成 }k\text{ 个 leaf vectors}.
\]

在本文中，我们也会把
\[
\bigoplus_{r=0}^{k-1}S_r
\]
与 \(S\) 视为同一个分组操作；前者强调它由 \(k\) 个选择算子组成，后者强调整体映射。

\begin{remark}
    
    这里的 \(S_r\) 不是一个需要实际在密文上执行的算法步骤，而是用来描述
\(\Split\) 后第 \(r\) 个 leaf 再经过 leaf \(\CtwoS_n\) 的明文坐标效果。

确切地说，我们有如下交换关系：
\[
\boxed{
\CtwoS_n\circ \Split_r
=
S_r\circ \CtwoS_N.
}
\]

证明如下。设大环明文多项式为
\[
P(X)=\sum_{i=0}^{N-1}a_iX^i,
\]
系数向量为
\[
p=(a_0,a_1,\ldots,a_{N-1})^T.
\]
大环 canonical slots 为
\[
z=F_Np.
\]
因此
\[
\CtwoS_N(z)=F_N^{-1}z=p.
\]

另一方面，coefficient split 给出
\[
P(X)=\sum_{r=0}^{k-1}X^rP_r(X^k),
\]
其中
\[
P_r(Y)=\sum_{j=0}^{n-1}a_{r+jk}Y^j.
\]
所以第 \(r\) 个 leaf 多项式 \(P_r\) 的系数向量是
\[
S_rp=(a_r,a_{r+k},a_{r+2k},\ldots,a_{r+(n-1)k})^T.
\]

第 \(r\) 个 leaf 的 canonical slots 是
\[
F_nS_rp.
\]
对它做 leaf \(\CtwoS_n\)，得到
\[
\CtwoS_n(F_nS_rp)
=
F_n^{-1}F_nS_rp
=
S_rp.
\]

而 direct 路径中，
\[
S_r\circ \CtwoS_N(F_Np)
=
S_rF_N^{-1}F_Np
=
S_rp.
\]

因此
\[
\CtwoS_n\circ \Split_r(F_Np)
=
S_r\circ\CtwoS_N(F_Np),
\]
即
\[
\boxed{
\CtwoS_n\circ \Split_r
=
S_r\circ \CtwoS_N.
}
\]

把所有 \(r=0,\ldots,k-1\) 合起来，得到
\[
\boxed{
\left(
\bigoplus_{r=0}^{k-1}\CtwoS_n
\right)
\circ
\Split
=
\left(
\bigoplus_{r=0}^{k-1}S_r
\right)
\circ
\CtwoS_N.
}
\]

\end{remark}
\section{S2C 与 Merge 的重组}

上一节说明了 EvalMod 与 coefficient 分组的交换性。设大环 \(\CtwoS_N\) 之后得到的 coefficient-packed vector 为
\[
x=(x_0,x_1,\ldots,x_{N-1}).
\]

\begin{remark}
    
    x 是C2S之后的slot 内容；它等于C2S之前明文多项式P（X）的系数向量。

    所以它有双重描述：

\begin{enumerate}
    \item 从 C2S 后的密文 slot 语义看，$x$ 是 slots 中的值；
    \item 从 C2S 前的多项式 $P(X)$ 看，$x$ 是 $P(X)$ 的 coefficients。
\end{enumerate}

这就是 “coefficient-packed vector” 的意思

\end{remark}
大环 EvalMod 的理想抽象是
\[
\EvalMod_N=f^{\oplus N},
\]
因此经过大环 EvalMod 后得到
\[
y=f^{\oplus N}(x)
=
(f(x_0),f(x_1),\ldots,f(x_{N-1})).
\]

这里 \(y\) 仍然是 coefficient-packed vector，只是每个 coordinate 已经完成了模约减近似。换句话说，如果
\[
x_i=m_i+q_0t_i+\varepsilon_i,
\]
那么
\[
y_i=f(x_i)\approx m_i.
\]

在 no-\(B_k\) 的 leaf 路线中，先通过 \(S\) 将大环 coefficient-packed vector 按 residue class 分成
\[
Sx=(S_0x,\ldots,S_{k-1}x).
\]
由于 EvalMod 与该分组操作交换，leaf EvalMod 后得到的 \(k\) 个输出向量等价于
\[
S f^{\oplus N}(x)
=
(S_0y,S_1y,\ldots,S_{k-1}y),
\]
其中
\[
y=f^{\oplus N}(x).
\]


或者表达为在 no-\(B_k\) 的 leaf 路线中，上一节的交换性给出：leaf EvalMod 后的 \(k\) 个输出向量分别是
\[
S_0y,\ S_1y,\ \ldots,\ S_{k-1}y,
\]
也就是
\[
(S_0 f^{\oplus N}(x),\ S_1 f^{\oplus N}(x),\ \ldots,\ S_{k-1} f^{\oplus N}(x)).
\]

现在需要说明的是：leaf \(\StwoC_n\) 和 \(\Merge\) 会把这些 reduced coefficient 子向量重组回大环 slot 表示。
\subsection{理想重组}

上一节中，我们已经把 EvalMod 后的结果记为
\[
y=f^{\oplus N}(x).
\]

这里需要再次强调 \(y\) 的语义。  
它不是 canonical slot vector，而是 EvalMod 后的 coefficient-packed vector。也就是说，
\[
y=(y_0,y_1,\ldots,y_{N-1})^T
\]
中的每个 \(y_i\) 应该被理解为某个大环多项式的第 \(i\) 个系数。

如果 C2S 前的大环多项式系数满足
\[
x_i=m_i+q_0t_i+\varepsilon_i,
\]
那么经过 EvalMod 后，
\[
y_i=f(x_i)\approx m_i.
\]
因此 \(y\) 可以理解为模约减后的大环系数向量。

换句话说，如果最终要恢复的大环多项式记为
\[
Y(X),
\]
那么理想情况下它应该是
\[
Y(X)=\sum_{i=0}^{N-1}y_iX^i.
\]

现在我们要说明：leaf \(\StwoC_n\) 加上 \(\Merge\) 正好会把这些 \(y_i\) 重新组成这个大环多项式，并最终得到它的 canonical slots。

\subsubsection*{第一步：把大环系数向量按 residue class 分成 leaf 系数向量}

因为
\[
N=kn,
\]
每个下标 \(i=0,\ldots,N-1\) 都可以唯一写成（大环的系数下标按照模 k 的余数r分组。）
\[
i=r+jk,
\]
其中
\[
0\le r<k,\qquad 0\le j<n.
\]

因此，大环系数向量
\[
y=(y_0,y_1,\ldots,y_{N-1})^T
\]
可以按 residue class 分成 \(k\) 组。

第 \(r\) 组是 （注：$S_r:\mathbb{C}^N\to\mathbb{C}^n$）
\[
S_ry
=
(y_r,y_{r+k},y_{r+2k},\ldots,y_{r+(n-1)k})^T.
\]

这个向量正好可以看成第 \(r\) 个小环多项式的系数向量。定义
\[
Y_r(Y)
=
\sum_{j=0}^{n-1}y_{r+jk}Y^j.
\]

于是
\[
S_ry
\]
就是 \(Y_r(Y)\) 的系数向量。

例如当 \(k=2\) 时，
\[
S_0y=(y_0,y_2,y_4,\ldots)^T,
\]
对应小环多项式
\[
Y_0(Y)=y_0+y_2Y+y_4Y^2+\cdots,
\]
而
\[
S_1y=(y_1,y_3,y_5,\ldots)^T,
\]
对应
\[
Y_1(Y)=y_1+y_3Y+y_5Y^2+\cdots.
\]

\subsubsection*{第二步：leaf \(\StwoC_n\) 把每个 leaf 的系数向量转成小环 slots}

现在考虑第 \(r\) 个 leaf。

它的 coefficient-packed vector 是
\[
S_ry.
\]

也就是说，当前 leaf 的 slots 中装的是小环多项式 \(Y_r(Y)\) 的系数：
\[
(y_r,y_{r+k},y_{r+2k},\ldots).
\]

leaf \(\StwoC_n\) 的作用是：把这些 slots 中的值当作小环多项式的系数，并转回小环 canonical slot 表示。

在理想模型下，
\[
\StwoC_n=F_n.
\]

所以第 \(r\) 个 leaf 经过 \(\StwoC_n\) 后得到
\[
F_nS_ry.
\]

这就是小环多项式 \(Y_r(Y)\) 的 canonical slots：
\[
F_nS_ry
=
\sigma_n(Y_r).
\]

所有 leaf 合起来，就得到
\[
(F_nS_0y,\ F_nS_1y,\ \ldots,\ F_nS_{k-1}y).
\]

这表示我们现在有 \(k\) 个小环密文，它们分别加密小环多项式
\[
Y_0(Y),Y_1(Y),\ldots,Y_{k-1}(Y)
\]
的 canonical slot 表示。

\subsubsection*{第三步：Merge 把这些小环多项式拼回大环多项式}

现在 \(\Merge\) 的理想语义是，把 \(k\) 个小环多项式
\[
Y_0(Y),Y_1(Y),\ldots,Y_{k-1}(Y)
\]
按照 coefficient split 的逆过程拼回一个大环多项式。

也就是说，定义
\[
Y(X)
=
\sum_{r=0}^{k-1}X^rY_r(X^k).
\]

这正是 coefficient split
\[
P(X)=\sum_{r=0}^{k-1}X^rP_r(X^k)
\]
的逆过程。

把
\[
Y_r(Y)=\sum_{j=0}^{n-1}y_{r+jk}Y^j
\]
代入上式，得到
\[
Y(X)
=
\sum_{r=0}^{k-1}X^r
\left(
\sum_{j=0}^{n-1}y_{r+jk}(X^k)^j
\right).
\]

因为
\[
(X^k)^j=X^{jk},
\]
所以
\[
Y(X)
=
\sum_{r=0}^{k-1}
\sum_{j=0}^{n-1}
y_{r+jk}X^rX^{jk}.
\]

也就是
\[
Y(X)
=
\sum_{r=0}^{k-1}
\sum_{j=0}^{n-1}
y_{r+jk}X^{r+jk}.
\]

由于每个下标 \(i=0,\ldots,N-1\) 都唯一地写成
\[
i=r+jk,
\]
所以上式正好等于
\[
Y(X)=\sum_{i=0}^{N-1}y_iX^i.
\]

这说明：\(\Merge\) 后得到的大环多项式，其系数向量正是
\[
y.
\]

\subsubsection*{第四步：回到大环 canonical slots}

前面说过，\(y\) 是一个大环系数向量。对应的大环多项式是
\[
Y(X)=\sum_{i=0}^{N-1}y_iX^i.
\]

大环 canonical slot vector 是
\[
\sigma_N(Y)=F_Ny.
\]

理想情况下，
\[
\StwoC_N=F_N.
\]

因此，大环 direct S2C 对 coefficient-packed vector \(y\) 的输出是
\[
\StwoC_N(y)=F_Ny.
\]

而 leaf 路线中，先对每个 \(S_ry\) 做 leaf \(\StwoC_n\)，再 \(\Merge\)，得到的也是同一个大环多项式
\[
Y(X)=\sum_{i=0}^{N-1}y_iX^i
\]
的 canonical slots。

因此，理想情况下有交换关系：
\[
\boxed{
\Merge
\circ
\left(
\bigoplus_{r=0}^{k-1}\StwoC_n
\right)
\circ
\left(
\bigoplus_{r=0}^{k-1}S_r
\right)
=
\StwoC_N.
}
\]

这个等式的含义是：

\[
\text{先把 coefficient-packed vector 按 residue class 分给各个 leaf，}
\]
\[
\text{每个 leaf 做小环 S2C，最后 Merge，}
\]
等价于
\[
\text{直接对整个大环 coefficient-packed vector 做大环 S2C。}
\]

\subsubsection*{把 EvalMod 的输出代回去}

因为 EvalMod 后的 coefficient-packed vector 是
\[
y=f^{\oplus N}(x),
\]
所以 leaf 路线中进入 leaf \(\StwoC_n\) 的 \(k\) 个向量是
\[
S_0y,\ldots,S_{k-1}y.
\]

也就是
\[
S_0 f^{\oplus N}(x),\ldots,S_{k-1} f^{\oplus N}(x).
\]

因此有
\[
\Merge
\circ
\left(
\bigoplus_{r=0}^{k-1}\StwoC_n
\right)
\left(
S_0 f^{\oplus N}(x),
\ldots,
S_{k-1} f^{\oplus N}(x)
\right)
=
\StwoC_N\left(f^{\oplus N}(x)\right).
\]

这正是说明：EvalMod 后的 leaf outputs 可以通过 leaf S2C 和 Merge 正确重组为 direct 大环 S2C 的输出。

\section{主定理：no-\(B_k\) 的 bootstrapping core 分解}

现在合成前面的三个部分。

\begin{definition}[理想 bootstrapping core]
定义大环理想 bootstrapping core：
\[
\BTS_N
:=
\StwoC_N\circ\EvalMod_N\circ\CtwoS_N.
\]
定义小环 leaf bootstrapping core：
\[
\BTS_n
:=
\StwoC_n\circ\EvalMod_n\circ\CtwoS_n.
\]
\end{definition}

\begin{theorem}[理想 no-\(B_k\) 分解]
在理想 full complex embedding 模型下，若：
\begin{enumerate}[label=(\arabic*)]
    \item \(\Split_k\) 是 coefficient-domain residue split；
    \item \(\Merge_k\) 是其逆向重组；
    \item \(\CtwoS_N,\CtwoS_n\) 分别为 \(F_N^{-1},F_n^{-1}\)；
    \item \(\StwoC_N,\StwoC_n\) 分别为 \(F_N,F_n\)；
    \item \(\EvalMod_N,\EvalMod_n\) 是同一个逐坐标标量函数 \(f\)；
\end{enumerate}
则
\[
\boxed{
\Merge_k
\circ
\left(
\bigoplus_{r=0}^{k-1}\BTS_n
\right)
\circ
\Split_k
=
\BTS_N.
}
\]
\end{theorem}

\begin{proof}
从右向左跟踪一个大环 slot input：
\[
z=F_Np.
\]

第一步，由 C2S-Split 交换图：
\[
\left(\bigoplus_r \CtwoS_n\right)\circ\Split_k(z)
=
\left(\bigoplus_r S_r\right)\circ\CtwoS_N(z)
=
(S_0p,S_1p,\ldots,S_{k-1}p).
\]

第二步，leaf EvalMod 分别作用：
\[
\left(\bigoplus_r \EvalMod_n\right)(S_0p,\ldots,S_{k-1}p)
=
(f(S_0p),\ldots,f(S_{k-1}p)).
\]
由于 \(f\) 逐坐标作用，
\[
f(S_rp)=S_rf(p).
\]
因此得到
\[
(S_0f(p),S_1f(p),\ldots,S_{k-1}f(p)).
\]

第三步，leaf S2C 与 Merge 重组：
\[
\Merge_k
\circ
\left(\bigoplus_r\StwoC_n\right)
(S_0f(p),\ldots,S_{k-1}f(p))
=
F_Nf(p).
\]

而 direct bootstrapping core 为
\[
\BTS_N(z)
=
\StwoC_N\circ\EvalMod_N\circ\CtwoS_N(F_Np)
=
F_Nf(p).
\]
两者相等。
\end{proof}

\begin{corollary}[带 normalization 的实现形式]
真实实现中若线性重组部分带全局 normalization factor \(\gamma_k\)，则可写为
\[
\boxed{
\Merge_k^{impl}
\circ
\left(
\bigoplus_{r=0}^{k-1}\BTS_n^{impl}
\right)
\circ
\Split_k^{impl}
\approx
\gamma_k \BTS_N^{impl}.
}
\]
因此归一化后有
\[
\boxed{
\gamma_k^{-1}
\Merge_k^{impl}
\circ
\left(
\bigoplus_{r=0}^{k-1}\BTS_n^{impl}
\right)
\circ
\Split_k^{impl}
\approx
\BTS_N^{impl}.
}
\]
\end{corollary}

\section{为什么当前路线不需要 \(B_k\)}

本节解释 \(B_k\) 的数学含义，以及为什么它不是 no-\(B_k\) 路线所需的操作。

\subsection{\(B_k\) 的含义}

设
\[
P(X)=\sum_{r=0}^{k-1}X^rP_r(X^k).
\]
令 \(\alpha\) 是大环 root，并令
\[
\theta=\alpha^k.
\]
则
\[
P(\alpha)
=
\sum_{r=0}^{k-1}\alpha^rP_r(\theta).
\]

因此，如果已知 leaf canonical slots
\[
(P_0(\theta),P_1(\theta),\ldots,P_{k-1}(\theta)),
\]
想恢复对应的大环 canonical slots
\[
P(\alpha),
\]
就需要一个局部 Vandermonde 型矩阵：
\[
B_{k,\theta}
=
\left(\alpha_{\theta,\ell}^{\,r}\right)_{\ell,r}.
\]

对于 \(k=2\)，这就是
\[
\begin{pmatrix}
P(\alpha)\\
P(-\alpha)
\end{pmatrix}
=
\begin{pmatrix}
1 & \alpha\\
1 & -\alpha
\end{pmatrix}
\begin{pmatrix}
A(\theta)\\
B(\theta)
\end{pmatrix}.
\]

所以 \(B_k\) 的作用是：
\[
\boxed{
\text{从 leaf canonical slots 恢复 big-ring canonical slots。}
}
\]

\subsection{no-\(B_k\) 路线的坐标流}

当前 no-\(B_k\) 路线是：
\[
\Split
\longrightarrow
\bigoplus_r\CtwoS_n
\longrightarrow
\bigoplus_r\EvalMod_n
\longrightarrow
\bigoplus_r\StwoC_n
\longrightarrow
\Merge.
\]

其坐标流为：
\[
\text{big slots}
\longrightarrow
\text{leaf canonical slots}
\longrightarrow
\text{leaf coefficient-packed slots}
\longrightarrow
\text{reduced leaf coefficient-packed slots}
\longrightarrow
\text{leaf canonical slots}
\longrightarrow
\text{big slots}.
\]

关键是：经过 leafC2S 后，数据已经不再处于 leaf canonical slot 坐标
\[
(P_r(\theta))_r,
\]
而是处于 coefficient-packed 坐标：
\[
(a_r,a_{r+k},a_{r+2k},\ldots).
\]

EvalMod 需要处理的正是这些 coefficient-like values：
\[
a_i=m_i+q_0t_i.
\]

因此 no-\(B_k\) 路线不需要先恢复 big-ring canonical slots。
它只需要让 leafC2S 取得 coefficient 子序列，然后对这些子序列逐坐标 EvalMod。

\subsection{为什么 \(B_k\) 不应插在 leafC2S 之后}

\(B_k\) 的公式
\[
P(\alpha)
=
\sum_r \alpha^rP_r(\alpha^k)
\]
适用于 leaf canonical slots
\[
P_r(\theta).
\]
但是 leafC2S 之后，slot 中装的是
\[
\text{coefficients of }P_r,
\]
而不是
\[
P_r(\theta).
\]

因此把裸 \(B_k\) 插在
\[
leafC2S \longrightarrow \EvalMod
\]
之间，在坐标语义上并不自然。它试图把 coefficient-packed slots 当成 canonical evaluation slots 来 lift，这会造成坐标系不匹配。

\begin{remark}
这也解释了为什么 \(B_k\) 在裸多项式求值坐标下数学正确，但不一定适用于 Lattigo 的 leafC2S 输出。
裸 \(B_k\) 解决的是 canonical slot lift 问题；no-\(B_k\) 路线解决的是 coefficient-subvector-wise EvalMod 问题。
\end{remark}

\section{真实实现中的布局与 scale 问题}

前面的推导使用了理想矩阵 \(F_N,F_n\)。
真实实现中还需要考虑若干 convention。

\subsection{Coefficient layout}

理论中我们写
\[
(E/O)\circ\CtwoS_N.
\]
但真实实现的 \(\CtwoS_N\) 输出未必按自然顺序
\[
(a_0,a_1,a_2,a_3,\ldots)
\]
排列。

它可能输出 blocked layout，例如
\[
(a_0,a_2,a_4,\ldots\mid a_1,a_3,a_5,\ldots),
\]
或者经过 bit-reversal、permutation、conjugation convention 等。

因此在实现层面，交换图更准确地写为
\[
\left(\bigoplus_r \CtwoS_n^{impl}\right)\circ\Split_k^{impl}
\approx
\Pi_{\mathrm{layout}}
\circ
\left(\bigoplus_r S_r\right)
\circ
\CtwoS_N^{impl},
\]
其中 \(\Pi_{\mathrm{layout}}\) 是实现相关的 permutation 或 block layout 变换。

\subsection{Real/Imag split}

CKKS bootstrapping 实现中 \(\CtwoS\) 通常输出两个 ciphertext：
\[
ct_{\mathrm{real}},\quad ct_{\mathrm{imag}}.
\]
理论上可以把它理解为两个独立的线性分支：
\[
\CtwoS_N^{real},
\qquad
\CtwoS_N^{imag}.
\]
no-\(B_k\) 的交换图需要分别在 real 和 imag 分支上成立。

\subsection{Scale 与 level}

真实实现中还要满足：
\begin{enumerate}[label=(\arabic*)]
    \item top 与 leaf 使用相同或兼容的 \(Q/P\) modulus chain；
    \item EvalMod 输入 scale 对齐；
    \item leaf EvalMod 的 approximation parameters 与 top 一致或等价；
    \item leaf S2C 的输入 level 满足实现要求；
    \item Split/Merge key switching 噪声可控；
    \item normalization factor \(\gamma_k\) 可被全局吸收。
\end{enumerate}

这些条件不改变理想分解定理，但决定真实系统中误差和 level 是否可接受。

\section{当前理论成果总结}

目前已经得到的理论认识可以总结如下：

\begin{enumerate}[label=(\arabic*)]
    \item Coefficient split 是
    \[
    P(X)=\sum_{r=0}^{k-1}X^rP_r(X^k),
    \]
    它不是 slot-preserving split。
    
    \item 对于 \(k=2\)，coefficient split 在 slot 视角下是
    \[
    (P(\alpha),P(-\alpha))
    \mapsto
    \left(
    \frac{P(\alpha)+P(-\alpha)}{2},
    \frac{P(\alpha)-P(-\alpha)}{2\alpha}
    \right).
    \]
    
    \item 若要 slot-preserving split，需要额外构造
    \[
    Q^\pm(Y)=A(Y)\pm h(Y)B(Y),
    \]
    其中
    \[
    h(\theta)=\alpha,\qquad h(Y)^2\equiv Y.
    \]
    
    \item 当前 no-\(B_k\) 路线使用 coefficient split，而不是 slot-preserving split。
    
    \item 理想模型下，C2S 与 coefficient split 满足交换图：
    \[
    \left(\bigoplus_r\CtwoS_n\right)\circ\Split_k
    =
    \left(\bigoplus_rS_r\right)\circ\CtwoS_N.
    \]
    
    \item 如果 EvalMod 是同一个逐坐标标量函数 \(f\)，则
    \[
    \left(\bigoplus_r\EvalMod_n\right)
    \circ
    \left(\bigoplus_rS_r\right)
    =
    \left(\bigoplus_rS_r\right)
    \circ
    \EvalMod_N.
    \]
    
    \item leaf S2C 与 Merge 可以将 reduced coefficient 子序列重组回大环 slot。
    
    \item 因此理想模型下：
    \[
    \Merge_k
    \circ
    \left(
    \bigoplus_r\BTS_n
    \right)
    \circ
    \Split_k
    =
    \BTS_N.
    \]
    
    \item 真实实现中可能出现全局 normalization factor：
    \[
    \Merge_k^{impl}
    \circ
    \left(
    \bigoplus_r\BTS_n^{impl}
    \right)
    \circ
    \Split_k^{impl}
    \approx
    \gamma_k\BTS_N^{impl}.
    \]
    
    \item \(B_k\) 是 leaf canonical slots 到 big canonical slots 的 lift 矩阵；no-\(B_k\) 路线在 leafC2S 后进入 coefficient-packed coordinates，因此不需要显式 \(B_k\)。
\end{enumerate}

\section{尚未完全解决的问题}

虽然理想理论已经给出清晰框架，但仍有若干问题需要继续研究：

\begin{enumerate}[label=(\arabic*)]
    \item \(\gamma_k\) 的精确来源是什么？
    
    当前猜测一般可能有
    \[
    \gamma_k=k.
    \]
    需要从 ring switching、trace/projection、Merge convention 或 S2C normalization 中推导。
    
    \item 对任意 \(k\) 的 layout 如何描述？
    
    对 \(k=2\)，可以用 even/odd 或 blocked layout 描述。
    对 \(k=4,8,\ldots\)，需要明确 leaf ordering、bit-reversal、residue class layout。
    
    \item Lattigo 的 C2S 实际矩阵与理想 \(F_N^{-1}\) 的关系是什么？
    
    理论中用 \(F_N^{-1}\)，实现中是 factorized DFT、permutation、real/imag split 和 scale 的组合。
    
    \item top 与 leaf EvalMod 参数何时完全等价？
    
    需要比较：
    \[
    q_0,\quad \log\text{MessageRatio},\quad K,\quad degree,\quad double\ angle,\quad scale.
    \]
    
    \item 误差如何累积？
    
    需要建立误差模型：
    \[
    \epsilon_{\mathrm{total}}
    =
    \epsilon_{\mathrm{split}}
    +
    \epsilon_{\mathrm{C2S}}
    +
    \epsilon_{\mathrm{EvalMod}}
    +
    \epsilon_{\mathrm{S2C}}
    +
    \epsilon_{\mathrm{merge}}
    +
    \epsilon_{\mathrm{keyswitch}}.
    \]
    
    \item level 消耗是否真的优于或不劣于 direct bootstrapping？
    
    如果 no-\(B_k\) 路线不再需要为 \(B_k,B_k^{-1}\) 预留 level，则输出 level 可能可以改善。
    
    \item 成本模型如何建立？
    
    需要比较：
    \[
    T_{\mathrm{direct}}
    \]
    与
    \[
    T_{\mathrm{factored}}
    =
    T_{\mathrm{Split}}
    +
    \max_r T_{\mathrm{leafBTS}}
    +
    T_{\mathrm{Merge}}.
    \]
    
    还需要计入 key generation、ring switching keys、leaf bootstrapping keys、memory footprint 等。
\end{enumerate}

\section{可能的后续理论方向}

\subsection{任意 \(k\) 的完整 slot-level 推导}

需要把 \(k=2\) 的 pairwise 推导推广到 \(k\)-fiber：
\[
\alpha_{\theta,0},\ldots,\alpha_{\theta,k-1},
\qquad
\alpha_{\theta,\ell}^k=\theta.
\]
Coefficient split 给出
\[
P(X)=\sum_{r=0}^{k-1}X^rP_r(X^k).
\]
在 slot 层面：
\[
P(\alpha_{\theta,\ell})
=
\sum_{r=0}^{k-1}
\alpha_{\theta,\ell}^rP_r(\theta).
\]
局部矩阵是 Vandermonde：
\[
B_{k,\theta}
=
\left(\alpha_{\theta,\ell}^{\,r}\right)_{\ell,r}.
\]
但 no-\(B_k\) 路线不需要显式使用它，因为 C2S 之后进入 coefficient coordinates。

\subsection{Normalization factor 的理论化}

需要分析实现中的 Split/Merge 是否满足：
\[
\Merge\circ\Split=\Id
\]
还是
\[
\Merge\circ\Split=\gamma_k\Id.
\]
也可能 Split/Merge roundtrip 本身是 identity，但
\[
\Merge\circ(\oplus \StwoC_n)
\]
相对
\[
\StwoC_N
\]
带有 \(\gamma_k\)。

这一点需要进一步区分。

\subsection{与 slot-preserving split 的关系}

可以把 coefficient split 和 slot-preserving split 写成两个不同路线：

\[
p\mapsto(S_0p,\ldots,S_{k-1}p)
\]
与
\[
p\mapsto F_n^{-1}SF_Np.
\]
前者高效、适合 no-\(B_k\) bootstrapping decomposition；
后者更适合真正的 slot extraction / slot distribution，但通常是稠密线性变换。

\section{结论}

本文理论框架说明，no-\(B_k\) subring bootstrapping decomposition 的关键并不是“小环 slots 是否等于大环 slots 的子集”，而是以下交换结构：

\[
\boxed{
\left(\bigoplus_r\CtwoS_n\right)\circ\Split_k
=
\left(\bigoplus_rS_r\right)\circ\CtwoS_N.
}
\]

Coefficient split 在 slot 视角下确实会混合大环 slots；但随后 leafC2S 会把这些 leaf slots 转回对应的 coefficient 子序列。
因为 EvalMod 是逐坐标函数，只要 top 与 leaf 使用相同的 scalar approximation，则 EvalMod 可以在每个 leaf 上独立执行。
最后 leafS2C 与 Merge 将 reduced coefficient 子序列重组回大环 slots。

因此，在理想模型下有：
\[
\boxed{
\Merge_k
\circ
\left(
\bigoplus_{r=0}^{k-1}\BTS_n
\right)
\circ
\Split_k
=
\BTS_N.
}
\]

在真实实现中，更合理的形式是：
\[
\boxed{
\Merge_k^{impl}
\circ
\left(
\bigoplus_{r=0}^{k-1}\BTS_n^{impl}
\right)
\circ
\Split_k^{impl}
\approx
\gamma_k\BTS_N^{impl}.
}
\]

这里 \(\gamma_k\) 是实现相关的全局 normalization factor。

最后，\(B_k\) 的作用是从 leaf canonical slots lift 到 big-ring canonical slots。
no-\(B_k\) 路线并不需要做这个 lift，因为它在 leafC2S 后进入了 coefficient-packed coordinates，EvalMod 正是在这些 coordinates 上逐坐标执行。
因此，显式 \(B_k\) 不是该分解理论的必要组成部分。

\section{核心交换图总结}

本文的核心观点可以浓缩为三个交换关系。它们分别对应 bootstrapping core 中的三段：
\[
\CtwoS
\quad\longrightarrow\quad
\EvalMod
\quad\longrightarrow\quad
\StwoC.
\]

设
\[
N=kn.
\]
令
\[
\Split_k
\]
表示 coefficient-domain split：
\[
P(X)=\sum_{r=0}^{k-1}X^rP_r(X^k),
\]
其中
\[
P_r(Y)=\sum_{j=0}^{n-1}a_{r+jk}Y^j.
\]

令
\[
S_r:\mathbb C^N\to\mathbb C^n
\]
为 residue-class 选择算子：
\[
S_r(x_0,\ldots,x_{N-1})
=
(x_r,x_{r+k},x_{r+2k},\ldots,x_{r+(n-1)k}).
\]

把所有 \(S_r\) 合并，定义
\[
S=
\bigoplus_{r=0}^{k-1}S_r.
\]

也就是
\[
Sx=(S_0x,S_1x,\ldots,S_{k-1}x).
\]

\subsection*{1. C2S 与 Split 的交换图}

第一个、也是最关键的交换图是：
\[
\boxed{
\left(
\bigoplus_{r=0}^{k-1}\CtwoS_n
\right)
\circ
\Split_k
=
S
\circ
\CtwoS_N.
}
\]

也可以逐 leaf 写成：
\[
\boxed{
\CtwoS_n\circ\Split_{k,r}
=
S_r\circ\CtwoS_N,
\qquad r=0,\ldots,k-1.
}
\]

其中 \(\Split_{k,r}\) 表示 split 后取第 \(r\) 个 leaf。

这个交换图的含义是：

\[
\text{先 coefficient split，再对每个 leaf 做 C2S}
\]
等价于
\[
\text{先做大环 C2S，再按 residue class 把 coefficients 分给各个 leaf。}
\]

也就是说，虽然 coefficient split 在 slot 视角下不是简单的 slot selection，但是
\[
\Split_k+\text{leaf }\CtwoS_n
\]
整体上正好实现了大环 C2S 输出的 coefficient 分组。

用图表示为：
\[
\begin{tikzcd}[column sep=large,row sep=large]
\text{big-ring canonical slots}
    \arrow[r,"\Split_k"]
    \arrow[d,"\CtwoS_N"']
&
\text{\(k\) leaf canonical slots}
    \arrow[d,"\bigoplus_r \CtwoS_n"]
\\
\text{big-ring coefficient-packed slots}
    \arrow[r,"S=\oplus_r S_r"']
&
\text{\(k\) leaf coefficient-packed slots}
\end{tikzcd}
\]

\subsection*{2. EvalMod 与 coefficient 分组的交换图}

设单坐标 EvalMod 标量函数为
\[
f.
\]
大环 EvalMod 抽象为
\[
\EvalMod_N=f^{\oplus N},
\]
小环 EvalMod 抽象为
\[
\EvalMod_n=f^{\oplus n}.
\]

由于 \(f\) 是逐坐标作用的，有：
\[
f^{\oplus n}(S_rx)=S_rf^{\oplus N}(x).
\]

因此：
\[
\boxed{
\left(
\bigoplus_{r=0}^{k-1}\EvalMod_n
\right)
\circ
S
=
S
\circ
\EvalMod_N.
}
\]

也就是：
\[
\boxed{
\left(
\bigoplus_{r=0}^{k-1} f^{\oplus n}
\right)
\circ
\left(
\bigoplus_{r=0}^{k-1}S_r
\right)
=
\left(
\bigoplus_{r=0}^{k-1}S_r
\right)
\circ
f^{\oplus N}.
}
\]

这个交换图的含义是：

\[
\text{先把 coefficients 分组，再在每个 leaf 上 EvalMod}
\]
等价于
\[
\text{先在大环上逐坐标 EvalMod，再把结果分组。}
\]

图示为：
\[
\begin{tikzcd}[column sep=large,row sep=large]
\text{big-ring coefficient-packed slots}
    \arrow[r,"S"]
    \arrow[d,"\EvalMod_N=f^{\oplus N}"']
&
\text{\(k\) leaf coefficient-packed slots}
    \arrow[d,"\bigoplus_r \EvalMod_n"]
\\
\text{reduced big-ring coefficient-packed slots}
    \arrow[r,"S"']
&
\text{\(k\) reduced leaf coefficient-packed slots}
\end{tikzcd}
\]

\subsection*{3. S2C 与 Merge 的交换图}

第三个交换图描述 leaf S2C 和 Merge 如何把分组后的 coefficient vectors 重组回大环 canonical slots：

\[
\boxed{
\Merge_k
\circ
\left(
\bigoplus_{r=0}^{k-1}\StwoC_n
\right)
\circ
S
=
\StwoC_N.
}
\]

这个交换图的含义是：

\[
\text{先把 reduced coefficients 分组，每组做 leaf S2C，再 Merge}
\]
等价于
\[
\text{直接对完整的大环 reduced coefficient vector 做大环 S2C。}
\]

图示为：
\[
\begin{tikzcd}[column sep=large,row sep=large]
\text{big-ring coefficient-packed slots}
    \arrow[r,"S"]
    \arrow[d,"\StwoC_N"']
&
\text{\(k\) leaf coefficient-packed slots}
    \arrow[d,"\bigoplus_r \StwoC_n"]
\\
\text{big-ring canonical slots}
&
\text{\(k\) leaf canonical slots}
    \arrow[l,"\Merge_k"']
\end{tikzcd}
\]

\subsection*{4. 合成后的 bootstrapping core 交换关系}

将以上三个交换图合成，得到 no-\(B_k\) bootstrapping core 的理想分解：
\[
\boxed{
\Merge_k
\circ
\left(
\bigoplus_{r=0}^{k-1}
\left(
\StwoC_n\circ\EvalMod_n\circ\CtwoS_n
\right)
\right)
\circ
\Split_k
=
\StwoC_N\circ\EvalMod_N\circ\CtwoS_N.
}
\]

也就是：
\[
\boxed{
\Merge_k
\circ
\left(
\bigoplus_{r=0}^{k-1}\BTS_n
\right)
\circ
\Split_k
=
\BTS_N.
}
\]

在真实实现中，如果 Split/Merge 或 S2C/Merge convention 带有全局 normalization factor \(\gamma_k\)，则更合理地写为
\[
\boxed{
\Merge_k^{impl}
\circ
\left(
\bigoplus_{r=0}^{k-1}\BTS_n^{impl}
\right)
\circ
\Split_k^{impl}
\approx
\gamma_k \BTS_N^{impl}.
}
\]

因此归一化后：
\[
\boxed{
\gamma_k^{-1}
\Merge_k^{impl}
\circ
\left(
\bigoplus_{r=0}^{k-1}\BTS_n^{impl}
\right)
\circ
\Split_k^{impl}
\approx
\BTS_N^{impl}.
}
\]

\subsection*{5. 解释}

以上交换图说明：

\[
\boxed{
no\text{-}B_k \text{ 路线不是在小环中保留原始大环 slots。}
}
\]

它真正利用的是：

\[
\boxed{
\Split_k+\text{leaf }\CtwoS_n
\text{ 等价于 }
\CtwoS_N
\text{ 后的 coefficient 分组。}
}
\]

因此，EvalMod 可以在各个 leaf 上独立执行，因为 EvalMod 是逐坐标函数。最后，leaf S2C 与 Merge 将 reduced coefficient 子向量重组回大环 canonical slots。

这就是 no-\(B_k\) bootstrapping decomposition 的核心理论结构。


把 Step 1, Step 2, Step 3 纵向串联，可得完整 bootstrapping core 分解图：

\[
\begin{tikzcd}[column sep=huge,row sep=large]
\mathbb{C}^{N}_{\mathrm{can}}
    \arrow[r,"\Split_k"]
    \arrow[d,"\CtwoS_N"']
&
\displaystyle\bigoplus_{r=0}^{k-1}\mathbb{C}^{n}_{\mathrm{can}}
    \arrow[d,"\displaystyle\bigoplus_{r=0}^{k-1}\CtwoS_n"]
\\
\mathbb{C}^{N}_{\mathrm{coeff}}
    \arrow[r,"S"]
    \arrow[d,"f^{\oplus N}"']
&
\displaystyle\bigoplus_{r=0}^{k-1}\mathbb{C}^{n}_{\mathrm{coeff}}
    \arrow[d,"\displaystyle\bigoplus_{r=0}^{k-1}f^{\oplus n}"]
\\
\mathbb{C}^{N}_{\mathrm{coeff}}
    \arrow[r,"S"]
    \arrow[d,"\StwoC_N"']
&
\displaystyle\bigoplus_{r=0}^{k-1}\mathbb{C}^{n}_{\mathrm{coeff}}
    \arrow[d,"\displaystyle\bigoplus_{r=0}^{k-1}\StwoC_n"]
\\
\mathbb{C}^{N}_{\mathrm{can}}
&
\displaystyle\bigoplus_{r=0}^{k-1}\mathbb{C}^{n}_{\mathrm{can}}
    \arrow[l,"\Merge_k"']
\end{tikzcd}
\]

\section{一句话总结}

线性操作可以FFT分解并行计算是显然的。我们证明了即使是非线性运算，只要它是对系数逐坐标作用的，也可以分解为并行计算。

\end{document}